% Part: proof-theory
% Chapter: normalization
% Section: normalization-thm

\documentclass[../../../include/open-logic-section]{subfiles}

\begin{document}

\olfileid{pt}{nor}{thm}

\olsection{في مبرهنة التطبيع}

نسمي~!!^a{proof} ذا صورة \emph{سوية} إذا خلا من مقاطع القطع. وهذا مكافئ
لألا يقبل تحويل تبديل ولا تحويل اختزال. و\emph{تطبيع}~!!a{proof} هو
نقله إلى برهان سوي، بإعمال تحويلات التبديل والاختزال حتى تنفد القطوع
كلها. وكون هذا ممكنًا دائمًا هو مضمون \emph{مبرهنة التطبيع}.

\begin{thm}\ollabel{thm:normalization}
كل~!!{proof}~$\delta$ في~\Log{N1i} قابل للتطبيع؛ أي إنه يتحول إلى
!!a{proof}~$\delta'$ خال من القطوع.
\end{thm}

\begin{proof}
  نسلك الاستقراء على رتبة القطع وطوله في~$\delta$. ونفرض أن كل~!!{proof}
  أدنى في رتبته القطعية، أو مساو فيها وأقصر طولًا قطعيًا، قابل للتطبيع.

  فإن كان طول القطع~$\OLMathNumeral{0}$ لم يكن قطع، وكان~$\delta$ سويًا من الأصل. فنفرض
  أن رتبة القطع وطوله في~$\delta$ كلاهما أكبر من~$\OLMathNumeral{0}$، ونختار مقطع قطع
  أقصى، علويًا، وفي أقصى اليمين. فإن زاد طوله على~$\OLMathNumeral{1}$ أعملنا تحويل تبديل؛
  وإن كان طوله~$\OLMathNumeral{1}$ أعملنا تحويل الاختزال الذي يقابله. فينشأ~!!{proof}
  رتبته القطعية هي نفسها وطوله القطعي أقصر، أو رتبته القطعية أدنى؛ وفي
  الحالين يعمل فرض الاستقراء.
\end{proof}

والطريقة التي سلكها البرهان تعين ترتيبًا مخصوصًا للتحويلات الواقعة على
!!a{proof} حتى تنتج~!!{proof} سويًا: نعالج في كل مرة مقطع قطع أقصى،
علويًا، وفي أقصى اليمين. فلا تثبت الحجة انتهاء التطبيع إلا لهذه الطريقة
المختارة. وأما دعوى انتهاء كل ترتيب اعتباطي لتحويلات التبديل والاختزال
إلى برهان سوي، فهي أقوى؛ وتسمى التطبيع القوي، ولا بد لها من مبرهنة أخرى
لا يفي بها هذا البرهان.

ويصاغ تعريف المقطع والقطع في~\Log{N2i} بغير عسر، وتنقل تحويلات التبديل
والاختزال من~\Log{N1i} إلى~\Log{N2i} مباشرة؛ فتصح مبرهنة التطبيع في
\Log{N2i} أيضًا.

\begin{thm}\ollabel{thm:normalization-N2i}
كل~!!{proof}~$\delta$ في~\Log{N2i} قابل للتطبيع؛ أي إنه يتحول إلى
!!a{proof}~$\delta'$ خال من القطوع.
\end{thm}

\end{document}
